\chapter{Less Known Facts}
\label{lnf}

\section{Hertz's forceless mechanics }
\label{lnf:hertz}

By the 1890s, several physicists --- Kirchhoff and Mach among them
--- had grown uneasy with the status of \emph{force} as a
fundamental concept of mechanics. Force is defined operationally
through $F=ma$, but that same equation is routinely used to
\emph{introduce} force, which makes the law read uncomfortably like
a definition dressed up as an empirical statement. In his
posthumously published \emph{Die Prinzipien der Mechanik in neuem
Zusammenhange dargestellt} (1894), Heinrich Hertz proposed to remove
the discomfort by removing force altogether: he set out to build the
whole of mechanics from just three primitive notions --- space,
time, and mass --- with no independent concept of force at
all~\cite{hertz}.

\subsection{The principle of the straightest path}

Consider $N$ mass points with positions $\mathbf x_i\in\mathbb R^3$
and masses $m_i$, and equip the $3N$-dimensional configuration space
$Q$ with the natural mass-weighted (kinetic-energy) metric,
\begin{equation}
  ds^2 \;=\; \sum_{i=1}^N m_i\, d\mathbf x_i\cdot d\mathbf x_i,
  \qquad
  T \;=\; \frac12\left(\frac{ds}{dt}\right)^{\!2}.
  \label{hertz-metric}
\end{equation}
Hertz's masses are never free in the naive sense: they are always
tied together by \emph{rigid, workless connections} (idealized
mechanical linkages that transmit force but do no work), which
restrict the accessible configurations to some submanifold
$\mathcal C\subset Q$. His fundamental postulate --- the
\emph{principle of the straightest path}
--- is then simply:

\begin{quote}
\emph{A system free of "force" moves along the straightest possible
path compatible with its constraints, at constant speed.}
\end{quote}

"Straightest possible" is to be read literally, in the geometric
sense: the representative point traces a \emph{geodesic} of the
metric~\eqref{hertz-metric} restricted to $\mathcal C$, i.e.\ a
curve whose acceleration (in the induced metric) is purely normal to
$\mathcal C$ and is entirely accounted for by the constraint
reaction. Formally, if the constraints are written $g_a(\mathbf
x)=0$, $a=1,\dots,k$, the equations of motion are
\begin{equation}
  m_i\,\ddot{\mathbf x}_i \;=\; \sum_a \lambda_a\,
  \frac{\partial g_a}{\partial \mathbf x_i},
  \label{hertz-eom}
\end{equation}
with the multipliers $\lambda_a$ fixed by demanding $g_a(\mathbf
x(t))\equiv 0$. There is nothing on the right-hand side of
\eqref{hertz-eom} except constraint reactions: no applied-force
term appears anywhere, by construction. Constant speed
($|ds/dt|=\mathrm{const}$, i.e.\ constant kinetic energy) follows
automatically, since ideal constraints do no work.

\subsection{Hidden masses, or where "force" goes}

Of course, ordinary bodies \emph{do} appear to exert forces on one
another --- gravity, contact forces, and so on all look like genuine
interactions between visible masses that are not rigidly connected
to anything. Hertz's resolution was to posit \emph{hidden masses}:
invisible mass points, rigidly linked to the visible ones (and, in
his more speculative moments, identified with a mechanical
aether), whose only role is to curve the constraint submanifold
$\mathcal C$ in just the right way. If the full configuration space
splits as $Q = Q_{\mathrm{vis}}\times Q_{\mathrm{hid}}$ and the
\emph{complete} system --- visible plus hidden --- moves along a
geodesic of $\mathcal C$, then projecting that geodesic motion down
onto $Q_{\mathrm{vis}}$ alone generically produces a curve that is
\emph{not} geodesic in the visible sector by itself: extra terms
appear in its equation of motion that are, from the point of view of
an observer with access only to $Q_{\mathrm{vis}}$, indistinguishable
from the effect of an applied force~\cite{hertz}.

This is a genuinely striking claim: on Hertz's account, force is not
absent from experience, only absent from the fundamental ontology.
Every observed force law would, in principle, be explained by some
specific (if permanently unobservable) arrangement of hidden mass
and rigid linkage. The same basic move --- explain an apparent force
as the shadow, on a submanifold, of pure geometric (geodesic) motion
in a larger space --- reappears twice in twentieth-century physics
in far more successful forms: in Kaluza--Klein theory, where the
electromagnetic force on a charged particle is recovered as geodesic
motion in a five-dimensional space, and, far more importantly, in
general relativity, where gravity itself --- which Newton had to
treat as a force acting at a distance --- becomes geodesic motion of
a free particle through curved four-dimensional spacetime, with no
"force" required at all. Hertz did not anticipate curved spacetime;
his geometry is Euclidean and the curvature is manufactured entirely
by hidden extra coordinates. But the structural resemblance to
Einstein's later, physically successful geometrization of gravity is
close enough that Hertz's book is still read today chiefly for that
reason, rather than as a working theory in its own
right.

\subsection{A second, more conservative geometrization}

It is worth noting that ordinary, force-\emph{having} mechanics can
also be written as a geodesic principle, without any hidden masses
at all, via Jacobi's form of the principle of least action. For a
conservative system with potential $V(\mathbf x)$ and fixed total
energy $E$, the true trajectories are exactly the geodesics of the
\emph{Jacobi metric}
\begin{equation}
  ds_J^2 \;=\; 2\bigl(E-V(\mathbf x)\bigr)\, ds^2 ,
  \label{jacobi-metric}
\end{equation}
a conformal rescaling of the mass metric~\eqref{hertz-metric} by
the local kinetic energy. So there are, in the end, two quite
different ways to eliminate force from the description of a system
governed by a potential: enlarge the configuration space with hidden
coordinates and keep the metric flat (Hertz), or keep the visible
configuration space as it is and let the metric itself become curved
and energy-dependent (Jacobi). Both reproduce exactly the same
observable trajectories. Neither is more "true" than the other; they
are, in Hertz's own vocabulary, different \emph{images} of the same
phenomena. That a concept as apparently basic as force can be made
to evaporate --- twice over, in two structurally different ways ---
without changing a single observable prediction is the real content
of this section: it shows "force" to be a feature of a chosen
representation of mechanics, not a fact about the world that every
correct representation must contain.

Hertz's own programme did not survive contact with its critics.
Almost every contemporary reviewer, Boltzmann included, objected
that positing a fresh, permanently unobservable mechanism of hidden
masses for every distinct force law (gravitation, electromagnetism,
\ldots) bought elegance at the price of testability, and the
research programme was essentially abandoned within a generation. It
survives today less as working physics than as an early, and
unusually explicit, demonstration that mechanics has more than one
internally consistent axiomatic foundation --- a lesson later put to
far more productive use.

\subsection{My Own Story}

In the middle of the eighties I studied a problem of quantization
\footnote{Later I describe this problem in more details}.
I have noticed that usual Hamiltonian
\[\mathcal{H}=\Sigma_1^n \frac{\vec p_j^2}{2 m_j} + U(\vec q_j) \quad j=1..n\]
is equivalent to the Hamiltonian defined on expanded phase space 
$(\vec p_j,p_{n+1},\vec q_j,q_{n+1})$:
\begin{equation}
   \label {lnf:hext}
    \mathcal{\tilde H}=\Sigma_1^n \frac{\vec p_j^2}{2 m_j} + U(\vec q_j) \frac{p_{n+1}^2}{2 p_0^2},
\end{equation}
where $p_0$ is arbitrary constant of dimensionality of momentum. 
Indeed, $p_{n+1}$ is integral of motion, fixing its value to be $p_0$ 
we get $ \mathcal{H}=\mathcal{\tilde H}$.

I knew nothing about Hertz's mechanics and called this 'effective mass 
mechanics'. Fortunately due to a number of reasons I did not published
my investigations and so escaped unintentional but clear plagiarism.


\section{Configuration Space:Deeper Insight}
\label{lnf:cs}

Classical mechanics is usually presented as the most settled part of
physics: three laws, a variational principle, and a semester of
worked examples. By the time a student reaches quantum theory or
general relativity, mechanics has already receded into the
background as a body of technique rather than a live subject. This
impression is not wrong, exactly, but it is incomplete. Even such simple concept as \emph{configuration space} is not trivial.

Let's discuss configuration space in some details. For $N$ point particles it
is 
\[\underbrace{\mathbb{R}^3 \times \mathbb{R}^3 \times \dots
\times\mathbb{R}^3}_{N \text{times}}\] 
which naturally isomorphic to
$\mathbb{R}^{3 N}$as vector space, topological space, and smooth manifold. But
they emphasize different structures and interpretations. For example natural
rotation is $SO(3) \subset SO(3N)$ not $SO(3N)$, kinetic energy has natural block
structure. 

Two-dimensional isotropic harmonic oscillator and two
identical, spatially distant, non-interacting one-dimensional oscillators have
exactly the same configuration space, $\mathbb{R}^2$, and, once the two
one-dimensional coordinates are laid side by side, exactly the same
Hamiltonian. This coincidence might seem to threaten the idea that
configuration space carries direct physical meaning: if the bare pair $(Q,H)$
cannot distinguish ``one particle moving in a plane'' from ``two particles
moving on separate, distant lines,'' in what sense does configuration space
represent physical positions at all? We examine this coincidence carefully,
show that it persists at every level of structure computable from $(Q,H)$ alone
--- classical trajectories, the symplectic structure,
--- and then identify precisely what additional data is needed
to tell the two systems apart:  an embedding of $Q$ into physical space
together with the action of the spatial isometry group, which is geometrically
realized in one case and merely accidental in the other;
and  the algebra of interactions that locality in physical
space picks out as natural, which governs how each system couples to a generic
environment. We package this as the \emph{physical dressing} of a bare
Hamiltonian system and argue that configuration space, taken on its own, is
best understood as a representation-dependent bookkeeping device rather than a
self-standing arena of physical positions: its physical reading requires this
further data as an independent postulate, not something read off the space
itself.

In elementary treatments of analytical mechanics, configuration space is often
introduced as though it simply \emph{were} the space of physical positions of a
system: a point $q\in Q$ is ``where the system is,'' and the Hamiltonian
$H:T^*Q\to\mathbb{R}$ tells us how that point moves. This picture is useful
and, for the simplest systems, harmless. But it invites a stronger reading on
which the pair $(Q,H)$ is taken to \emph{fix} the physical system being
described --- as if knowing the configuration manifold and the Hamiltonian
function on it were enough to know what kind of physical situation one is
dealing with.

A simple example shows this stronger reading to be false. Consider, on the one
hand, a single particle of mass $m$ moving in a plane under an isotropic
harmonic restoring force (the two-dimensional isotropic oscillator), and, on
the other hand, two particles of equal mass $m$, held very far apart, each
independently oscillating along its own line under a harmonic force of the same
strength. Written in the natural coordinates for each case, the two systems
have configuration space $Q$ and Hamiltonian $H$
\emph{identically}. There is no sense in which the bare mathematical data
$(Q,H)$ --- or, equivalently, the classical trajectories it generates,
--- differs between the two cases. Yet no one would say that a
single particle vibrating in a plane and two distant, independent oscillators
are ``the same physical system.'' Something is being smuggled in from outside
$(Q,H)$ whenever we say what the system \emph{is}.

Let us consider two systems:

\paragraph{System A (the planar oscillator).} A single particle of mass $m$
moves in a plane, subject to an isotropic harmonic potential centered on a
fixed equilibrium point.

\paragraph{System B (two distant oscillators).} Two particles, each of mass
$m$, sit at fixed, well-separated base points in physical space and each
oscillates, independently and without interaction, along its own fixed axis.

Any account of what
distinguishes the two systems must therefore appeal to something that is not a
function of $(Q,H)$ alone.
We identify  independent sources of physical structure that are not
visible in $(Q,H)$, each of which distinguishes System A from System B once
made explicit.

\subsection{Embedding in physical space and the status of the rotational symmetry}

The transformation $R_\theta$ identified above is, as a transformation of the
abstract manifold $Q=\mathbb{R}^2$, a symmetry of $H$ in \emph{both} systems.
But this bare fact does not tell us whether $R_\theta$ corresponds to anything
happening in physical space. To ask that question we must specify an embedding
of $Q$ into physical space $\Sigma$ (taken here to be $\mathbb{R}^3$, or a
plane within it) that says what each coordinate \emph{is}.

For System A, the natural embedding sends the single particle's configuration
to its actual position,
\[
\iota_A(q_1,q_2)=\mathbf{X}^{(0)}+q_1\hat{e}_1+q_2\hat{e}_2 \ \in\ \Sigma,
\]
and $R_\theta$ is exactly the pushforward, under $\iota_A$, of a genuine
rotation of physical space about the fixed point $\mathbf{X}^{(0)}$: rotating
the lab by angle $\theta$ carries the particle from $\iota_A(q_1,q_2)$ to
$\iota_A(R_\theta(q_1,q_2))$. The symmetry is \emph{geometric}: it is
implemented by an isometry of $\Sigma$, it corresponds to a real conserved
angular momentum exchangeable with other systems via torque, and it is what
makes $L=q_1p_2-q_2p_1$ deserve the name ``angular momentum'' rather than being
merely a conserved combination of phase-space coordinates.

For System B, the natural embedding sends the pair of independent
one-dimensional configurations to the positions of two separate particles,
\[
\iota_B(q_1,q_2)=\Big(\mathbf{X}_1^{(0)}+q_1\hat{n}_1,\ \
\mathbf{X}_2^{(0)}+q_2\hat{n}_2\Big)\ \in\ \Sigma\times\Sigma,
\]
with $\mathbf{X}_1^{(0)}\neq\mathbf{X}_2^{(0)}$ (this is what ``distant''
means). Here $R_\theta$ cannot be the pushforward of \emph{any} isometry of
$\Sigma$.

\begin{proposition}[Non-transfer of the rotational symmetry]
For $\mathbf{X}_1^{(0)}\neq\mathbf{X}_2^{(0)}$, there is no isometry $\rho$ of
    $\Sigma$ such that $\iota_B\circ R_\theta=(\rho,\rho)\circ\iota_B$ holds
    for all $\theta$ (other than $\theta\in\pi\mathbb{Z}$).
\end{proposition}

\begin{proof}
An isometry $\rho$ of physical space acts \emph{pointwise}: the image
    $\rho(\mathbf{y})$ of a point $\mathbf{y}\in\Sigma$ depends only on
    $\mathbf{y}$ itself. In particular, the would-be image of particle 2's
    position, $\rho\big(\mathbf{X}_2^{(0)}+q_2\hat n_2\big)$, can depend only
    on $q_2$ --- it cannot depend on $q_1$, since particle~1's displacement
    plays no role in specifying the point $\mathbf{X}_2^{(0)}+q_2\hat{n}_2$
    that $\rho$ is asked to act on. But the second component of
    $\iota_B(R_\theta(q_1,q_2))$ is
    $\mathbf{X}_2^{(0)}+(q_1\sin\theta+q_2\cos\theta)\hat n_2$, which for
    $\sin\theta\neq 0$ depends non-trivially on $q_1$. No pointwise map can
    reproduce this dependence, so no single isometry $\rho$ can serve as
    $\iota_B\circ R_\theta=(\rho,\rho)\circ\iota_B$ for such $\theta$.
\end{proof}

The obstruction is not a matter of fine-tuning
$\mathbf{X}_1^{(0)},\mathbf{X}_2^{(0)},\hat n_1,\hat n_2$; it holds however
these are chosen, and it fails only in the degenerate limit
$\mathbf{X}_1^{(0)}=\mathbf{X}_2^{(0)}$ --- precisely the limit in which System
B stops being ``two distant oscillators'' and coincides with System A. So
$R_\theta$ is, for System B, a symmetry of the Hamiltonian with no geometric
referent: it mixes the coordinates of two causally and spatially independent
objects in a way that no rotation, translation, or reflection of the physical
world could ever do. It is what is sometimes called a \emph{hidden},
\emph{accidental}, or purely \emph{dynamical} symmetry --- real as a statement
about the equations of motion, but carrying none of the physical content
(conservation via Noether's theorem applied to an actual spacetime symmetry,
exchangeability of angular momentum with other systems, appearance in a
Wigner-Eckart-type selection rule tied to real rotations) that the same-looking
symmetry carries in System A.

\begin{remark}
This is not special to the isotropic case $\omega_1=\omega_2$. If the two
    frequencies differ, $H=\frac{p_1^2+p_2^2}{2m}+\frac12
    m(\omega_1^2q_1^2+\omega_2^2q_2^2)$ still describes, term for term, both
    ``one particle in an anisotropic planar potential'' and ``two distant,
    differently-tuned oscillators,'' and the argument of Section~2 goes through
    unchanged. The isotropic case is simply the sharpest illustration, because
    there the enhanced symmetry make the geometric/accidental distinction most
    conspicuous.
\end{remark}

\subsection{Locality and the coupling to a generic environment}
\label{loc}

A second, complementary way to separate the two systems is to ask not about
their isolated dynamics but about how they couple to the rest of the world.
Isolation is always an idealization; a real oscillator sits in some environment
(radiation field, phonon bath, measuring apparatus), and it is a basic fact
about physical interactions that they are \emph{local}: an environmental degree
of freedom that is spatially near one object couples appreciably to that object
and negligibly to a second object located far away.

For System B this has an immediate and specific consequence. Because particle~1
and particle~2 occupy distant regions of space, generic environmental
influences factorize: to good approximation, particle~1 couples to a bath $B_1$
associated with its own neighborhood and particle~2 couples to an independent
bath $B_2$ associated with its neighborhood,
\begin{equation}
H_{SE}\ \approx\ q_1\otimes \hat B_1 \ +\ q_2\otimes \hat B_2 ,
\label{localcoupling}
\end{equation}
with $\hat B_1,\hat B_2$ statistically independent (uncorrelated) environmental
operators. This is simply the statement that two distant systems have
independent noise sources --- an instance of what is often called cluster
decomposition.

For System A, by contrast, there is only one particle, occupying one region of
space, and a generic local environmental influence acts at that single
location. Such an influence couples to some particular direction $\hat
n\cdot(q_1,q_2)$ in configuration space determined by the geometry of the
perturbation (e.g.\ the direction of a stray field at the particle's location),
not separately and independently to $q_1$ and $q_2$; arranging two independent
environments that couple exclusively along the $q_1$-axis and the $q_2$-axis
respectively, at one and the same spatial point, is not something locality
delivers for free, the way \eqref{localcoupling} is delivered for free in
System B --- it would need to be imposed by hand. There is consequently no
generic physical reason for a local environment to single out the particular
coordinate split $(q_1,q_2)$, as opposed to any rotated pair
$(q_1',q_2')=R_\theta(q_1,q_2)$.

\begin{remark}
This argument is offered as a physical heuristic, not a theorem: it says what
    is \emph{generic} given locality, not what is strictly forced by $(Q,H)$
    (which, after all, says nothing at all about any environment). That is
    exactly the point --- the distinction it draws on is additional structure,
    namely a specification of which couplings to the wider world are
    ``local''/``natural'' for the system in question.\footnote{ This is close in spirit
    to the observation, developed rigorously by Zanardi, Lidar, and Lloyd
    \cite{zll}, that a tensor-product (subsystem)
    decomposition of a quantum state space is not read off the Hilbert space or
    the Hamiltonian in isolation but is induced by the algebra of operationally
    accessible interactions and measurements. Here the analogous claim is made
    already at the level of which classical/environmental couplings are
    physically natural.}
\end{remark}

\paragraph{Normal modes.} The same moral appears, in a sense running the other
way, in the theory of coupled oscillators. A collection of masses connected by
springs, in general spatially local and physically coupled, has a Hamiltonian
that, after diagonalization in terms of normal-mode coordinates, takes exactly
the form of a set of \emph{decoupled} one-dimensional oscillators, $H=\sum_k
\big(P_k^2/2+\frac12\Omega_k^2 Q_k^2\big)$. If one only looked at this final,
diagonal $(Q,H)$, one might be tempted to read it as describing independent,
uncoupled particles. But the normal-mode coordinates $Q_k$ are collective
variables, each typically a linear combination of every original particle's
displacement; a perturbation local to one physical mass is, in normal-mode
space, generically \emph{non-local}, spreading across many or all of the $Q_k$.
This is the mirror image of subsection~\ref{loc}:
there, an abstractly ``local-looking''
coordinate split ($q_1,q_2$) failed to be the physically local one for System
A; here, an abstractly decoupled-looking Hamiltonian hides a physically
coupled, spatially local system. Both cases show that reading physical locality
or particle-independence directly off the form of $H$ in a given coordinate
system is unreliable without independent knowledge of the embedding.

